\begin{acknowledgement}
撰寫致謝時，請打開 \code{contents/front/acknowledgement.tex}，保留 \code{acknowledgement} 環境，再把這些說明換成自己的文字。模板會自動加入標題，內文依一般段落書寫即可，無需另外加上章節編號。

可以先回想研究過程中受到哪些協助，例如討論研究問題、取得資料，或協助檢查論文內容，再依序寫下想感謝的人與具體事情。完成後，請確認姓名、職稱與計畫名稱是否正確，也可以重新讀一次，看看是否有遺漏需要致謝的對象。

若不需要致謝頁，請到 \code{main.tex} 移除引用這個檔案的 \code{\input}，再重新編譯確認裝訂順序。若只是要更換標題，例如改為「謝詞」，則在 \code{config.yml} 的 \code{sections:} 下設定 \code{acknowledgement: "謝詞"} 即可。
\end{acknowledgement}
